\section{Considerações Finais}
\label{sec:consideracoes}

Este relatório documentou as seis questões do trabalho final sobre adaptação de
modelos de linguagem ao domínio público piauiense, das técnicas de treinamento
(pré-treino, SFT e QLoRA) às de compressão, recuperação e segurança (destilação, RAG
e \textit{guardrails}), sempre com avaliação estruturada antes e depois de cada
intervenção.

A Questão 01 confirmou sua hipótese de forma robusta: o pré-treino continuado sobre
71.431 documentos do DOMPI-2025 reduziu a perplexidade de 21,78 para 2,09 e elevou a
acurácia de previsão de tokens de 42,25\% para 82,03\%, sem \textit{overfitting}. O
\textit{benchmark} de 25 perguntas, respondido pelos dois modelos, mostrou que o
ganho vai além do estilo: o modelo treinado passou a acertar fatos exatos do corpus,
como CNPJs e CEPs de prefeituras, saindo de zero acertos exatos para seis.

A Questão 02 demonstrou o ajuste fino supervisionado com um pipeline de dados
inteiramente programático, que transformou 13.762 registros ruidosos do docentesDC
em 1.977 pares de instrução verificáveis por construção. O treino reduziu a
perplexidade sobre os tokens de resposta em 48,6\% e fez o modelo responder
corretamente fatos que antes desconhecia, com o episódio de \textit{overfitting} na
terceira época neutralizado pela seleção automática do melhor \textit{checkpoint}.

A Questão 03 evidenciou o argumento central da adaptação eficiente: com 0,44\% dos
parâmetros e 1,77~GB de memória, o QLoRA superou a perplexidade do ajuste fino
completo, ao custo de menor retenção factual no modelo de 0,5 bilhão. O experimento
adicional com o \texttt{Qwen2.5-1.5B} quantizado obteve o melhor resultado geral do
pós-treino (perplexidade 4,75), sugerindo que a combinação de escala com adaptação
eficiente é o melhor uso de um orçamento fixo de memória.

A Questão 04 aplicou destilação por \textit{Chain-of-Thought} de um professor
\texttt{Qwen2.5-3B} para um aluno \texttt{Qwen2.5-1.5B}, confirmando a transferência
de conhecimento nas métricas quantitativas (perplexidade 22,9\% menor), mas revelando,
pela avaliação semântica com \textit{LLM-as-Judge}, que menor perplexidade não implica
respostas mais confiáveis, um alerta contra conclusões baseadas apenas em
correspondência textual.

A Questão 05 construiu uma aplicação de RAG do tipo \textit{Standard} sobre o
docentesDC, avaliada pelo \textit{framework} Ragas em 30 perguntas. Ainda que as
pontuações absolutas tenham sido limitadas por lacunas do \textit{benchmark} e por uma
discrepância de metadados na avaliação, a comparação qualitativa demonstrou o valor da
recuperação: com contexto, o sistema deixa de alucinar e passa a citar o material real
ou a admitir honestamente a ausência do dado.

A Questão 06 adicionou uma camada de \textit{guardrails} ao modelo pós-treinado na
Questão 02 e mediu um grau de proteção elevado: 100\% das ameaças bloqueadas com
apenas 5\% de falsos positivos e 95\% de disponibilidade, endereçando na prática o
dilema entre utilidade e segurança.

Como trabalhos futuros, pretende-se integrar as peças desenvolvidas em um único
assistente para o domínio piauiense, combinando um modelo adaptado, recuperação de
contexto e camada de proteção; ampliar o dataset de destilação e avaliá-lo fora da
distribuição de treino; adotar recuperação hierárquica e saneamento do
\textit{benchmark} no RAG; e substituir os trilhos por \textit{regex} dos
\textit{guardrails} por classificadores semânticos.
